% ======================================================================
% Journal of Cheminformatics -- setup snippets for the Springer Nature
% LaTeX template (December 2024 package).
%
% These are THREE separate blocks. Each goes in a different place in
% sn-article.tex. Do not paste this file as one piece.
% ======================================================================


% ----------------------------------------------------------------------
% BLOCK 1 -- the \documentclass line
%
% Find the ONE line in sn-article.tex that starts with \documentclass
% and is NOT commented out (no %% in front). Replace it with this line.
% Keep all the commented %%\documentclass lines around it as they are.
% ----------------------------------------------------------------------

\documentclass[pdflatex,sn-basic,Numbered,referee,lineno]{sn-jnl}

%   pdflatex  -- the journal compiles submissions with pdfLaTeX
%   sn-basic  -- labelled "Chemistry Reference Style" in the template
%   Numbered  -- numbered citations [1], [2], ...
%   referee   -- double line spacing (required by the journal)
%   lineno    -- line numbers in the margin (required by the journal)
%
% If compilation fails, the log names the offending option. Remove that
% one option and recompile; do not remove them all at once.


% ----------------------------------------------------------------------
% BLOCK 2 -- the abstract
%
% Find \abstract{ in sn-article.tex and replace the whole braced block.
% Whole abstract: 350 words maximum. Not structured -- no Background /
% Methods / Results subheadings. The Scientific Contribution sits INSIDE
% the abstract, 3 sentences maximum.
% ----------------------------------------------------------------------

\abstract{[Aim, approach, and principal findings, in continuous prose.
No subheadings. 350 words maximum including the sentences below.]

\textbf{Scientific Contribution} [Sentence 1: what this work establishes
that was not established before.] [Sentence 2: how it differs from prior
work on this topic.] [Sentence 3: why that difference matters for how
models are validated. -- Maximum three sentences.]}


% ----------------------------------------------------------------------
% BLOCK 3 -- the back matter
%
% In sn-article.tex, find \backmatter. Delete everything from the line
% AFTER \backmatter down to (but not including) \begin{appendices}, or
% down to \bibliography{sn-bibliography} if there is no appendix.
% Then paste the block below in its place.
%
% Note: BMC journals put Acknowledgements INSIDE Declarations. The
% template's default puts it outside, and lists medical items
% (Ethics approval, Consent to participate) that this journal does not use.
% ----------------------------------------------------------------------

\backmatter

\section*{List of abbreviations}

AUROC, area under the receiver operating characteristic curve; BH,
Benjamini--Hochberg; DICT, drug-induced cardiotoxicity; DILI,
drug-induced liver injury; ECE, expected calibration error; MCC,
Matthews correlation coefficient; MCE, maximum calibration error.
% [Edit this list to match the abbreviations the final manuscript uses.]

\section*{Declarations}

\bmhead{Availability of data and materials}

[State where the datasets and the analysis code are deposited, with the
DOI. The journal requires source code for all software needed to reach
the conclusions.]

\bmhead{Competing interests}

The author declares no competing interests.

\bmhead{Funding}

[State funding sources, or: No funding was received for this work.]

\bmhead{Authors' contributions}

[Single author: state that the author designed the study, performed the
analyses and wrote the manuscript.]

\bmhead{Acknowledgements}

[Optional. If a large language model was used in preparing the
manuscript, Springer Nature policy requires that use to be disclosed
here -- name the tool and say what it was used for.]

% Authors' information -- optional, omit if not needed.
